\section{Introduction}

Scaling has enabled Large Language Models (LLMs) to achieve state-of-the-art performance on a range of Natural Language Processing (NLP) tasks~\citep{Wang2018GLUEAM, Wang2019SuperGLUEAS, Rajpurkar2016SQuAD1Q}. More importantly, new capabilities have emerged from LLMs as they are scaled to hundreds of billions of parameters~\citep{wei2022emergent}: in-context few-shot learning \citep{Brown2020LanguageMA} makes it possible for an LLM to perform well on a task it never trained on with only a handful of examples; Chain-of-Thought (CoT) prompting \citep{Wei2022ChainOT, step_by_step} demonstrates strong reasoning ability of LLMs across diverse tasks with or without few-shot examples; self-consistency \citep{Wang2022SelfConsistencyIC} further improves the performance via self-evaluating multiple reasoning paths.

Despite these incredible capabilities of models trained on large text corpus~\citep{Brown2020LanguageMA,Chowdhery2022PaLMSL}, fundamentally improving the model performances beyond few-shot baselines still requires finetuning on an extensive amount of \textit{high-quality supervised} datasets. FLAN~\citep{wei2021finetuned,chung2022scaling} and T0~\citep{sanh2021multitask} curated tens of benchmark NLP datasets to boost zero-shot task performances on unseen tasks; InstructGPT~\citep{ouyang2022training} crowd-sourced many human answers for diverse sets of text instructions to better align their model to human instructions. %; Minerva~\citep{todo} parsed the full ArXiv database carfully for relevant articles to excel on challenging competitive math and science datasets.  
% Arguably, the paradigm of training an LLM remains largely unchanged. % -- to improve its performance on challenging tasks, an extensive amount of supervised data is required. 
While significant efforts were committed on collecting high-quality supervised datasets, human brain, on the contrary, is capable of the metacognition process \citep{dunlosky2008metacognition}, where we can refine our own reasoning ability without external inputs. % We can come up with a random problem, try to solve it, and summarize what we learned, all inside our brain without external inputs. 
%Large performance gains from self-consistency~\citep{Wang2022SelfConsistencyIC} also hint at the reservoir of \textit{dark} knowledge~\citep{hinton2015distilling,korattikara2015bayesian} inside a LLM that can be extracted for improving itself.

In this paper, we study how an LLM is able to \textit{self-improve} its reasoning ability without supervised data. We show that using only
%a few prompt exemplars, and 
input sequences (without ground truth output sequences) from multiple NLP task datasets, a pre-trained LLM is able to improve performances for both in-domain and out-of-domain tasks. 
Our method is shown in Figure~\ref{fig:framework}: we first sample multiple predictions using few-shot Chain-of-Thought (CoT)~\citep{Wei2022ChainOT} as prompts, filter ``high-confidence'' predictions using majority voting~\citep{Wang2022SelfConsistencyIC}, and finally finetune the LLM on these high-confidence predictions. The resulting model shows improved reasoning in both greedy and multi-path evaluations.
%As shown in Figure \ref{fig:framework}: we first collect ``well-thought'' predictions (few-shot chain-of-thought prompting followed by self-consistency) using a pre-trained LLM, then fine-tune the LLM to let it directly output (without self-consistency) those ``well-thought'' predictions. 
We call the model fine-tuned in this way as \textbf{Language Model Self-Improved (LMSI)}. This is similar to how a human brain sometimes learns: given a question, think multiple times to derive different possible results, conclude on how the question should be solved, and then learn from or memorize its own solution.
We empirically verify our method using a pre-trained \bigmodel~LLM, where our method not only improves training task performances (74.4\%$\rightarrow$82.1\% on GSM8K, 78.2\%$\rightarrow$83.0\% on DROP, 90.0\%$\rightarrow$94.4\% on OpenBookQA, and 63.4\%$\rightarrow$67.9\% on ANLI-A3), but also enhances out-of-domain (OOD) test tasks (AQUA, 
% 35.8\%$\rightarrow$39.0\% on 
% 75.3\%$\rightarrow$77.8\% on 
StrategyQA, 
% 72.0\%$\rightarrow$81.8\% on 
MNLI), achieving state-of-the-art performances in many tasks without relying on supervised ground truth answers. Lastly, we conduct preliminary studies on self-generating additional input questions and few-shot CoT prompts, which could further reduce the amount of human effort required for model self-improving, and ablation studies on important hyperparameters of our approach. 
%To verify our method empirically, we use a pre-trained model with 118 layers and 540 billion parameters. For each input sequence, we apply few-shot chain-of-thought prompting and self-consistency to obtain the ``well-thought'' prediction using the pre-trained model. We then fine-tune the model on the well-thought prediction directly, and on all of the chain-of-thought sequences that lead to the well-thought predictions. 
%Our method improved the performance on GSM8K from 74\% to 82\% in terms of exact-match accuracy scores. We achieve state-of-the-art-level results on multiple benchmarks.
%The fact that our self-improving method does not require ground truth labels from training datasets makes large-scale fine-tuning more feasible, since collecting input sequences (questions) only requires less human effort. However, in our self-improving workflow, a small amount of human effort is still needed to: 1) collect input questions; 2) write in-context few-shot exemplars. To address these two problems and make our framework entirely \textit{unsupervisd}, we further investigate on how to generate a large amount of input questions and in-context few-shot exemplars, using the model itself. Early experiments show promising results.
We hope our simple approach and strong empirical results could encourage more future work by the community to investigate optimal performances of pretrained LLMs without additional human supervision.

Our contributions are summarized as follows:
\begin{itemize}
    \item We demonstrate that a large language model can self-improve by taking datasets without ground truth outputs, by leveraging CoT reasoning~\citep{Wei2022ChainOT} and self-consistency~\citep{Wang2022SelfConsistencyIC}, achieving competitive in-domain multi-task performances as well as out-of-domain generalization. We achieve state-of-the-art-level results on ARC, OpenBookQA, and ANLI datasets.
    \item We provide detailed ablation studies on training sample formatting and sampling temperature after fine-tuning, and identify critical design choices for most successful self-improvement by LLMs.
    \item We study two other approaches for self-improvements, where the model generates additional questions from finite input questions and generates few-shot CoT prompt templates itself. The latter achieves 74.2\% on GSM8K, which is the state-of-the-art \textit{zero-shot} performance, against 43.0\% by~\citet{step_by_step} or 70.1\% through its naive extension with~\citet{Wang2022SelfConsistencyIC}. 
   % \item We demonstrate that a model can improve itself by taking datasets without ground truth labels as input. This is achieved by fine-tuning on the model's own generated responses which are not necessarily always correct.
    %\item We show that it is possible to generate input questions or few-shot chain-of-thought exemplars. Using those generated questions and exemplars, our complete, end-to-end self-improving workflow is human-input free.
    %\item We achieve state-of-the-art-level results on multiple benchmark datasets: LIST THEM.
\end{itemize}

%The rest of this paper is organized as follows. Section \ref{sec:related_work} discusses related work. Section \ref{sec:method} lays out our method in detail. Section \ref{sec:setup} shows our setup for experiments. Section \ref{sec:results} shows our  results with ablation studies. Section \ref{sec:conclusions} concludes our work. The chain-of-thought prompts used in our work are included in Appendix \ref{app:appendix}.

\begin{figure*}[t]
% \subfigcapmargin=10pt
\centering
\includegraphics[width=1.0\textwidth]{Figures/method.pdf}
\vspace{0.2em}
\caption{Overview of our method. With Chain-of-Thought (CoT) examples as demonstration~\citep{Wei2022ChainOT}, the language model generates multiple CoT reasoning paths and answers (temperature $T>0$) for each question. The most consistent answer is selected by majority voting~\citep{Wang2022SelfConsistencyIC}. The ``high-confidence'' CoT reasoning paths that lead to the majority answer are augmented by mixed formats as the final training samples to be fed back to the model for fine-tuning.
% : (a) the training question is \textbf{preceded by CoT examples} and followed by a CoT reasoning path; (b) the training question is \textbf{preceded by standard prompting examples} and followed by a direct answer; (c) no examples are provided and ``Let's think step by step'' is used to \textbf{guide the model to generate CoT reasoning paths}; (d) no examples are provided and the model is trained to \textbf{directly predict the answer}.
}
\label{fig:framework}
% \vspace{-0.3cm}
% \vspace{-0.5em}
\end{figure*}